\documentclass[12pt]{article}
\usepackage[utf8]{inputenc}
\usepackage[T2A]{fontenc}
\usepackage[russian]{babel}
\usepackage{amsmath, amssymb, amsthm}

\begin{document}

\section*{Сопряжённые функции}

\subsection*{Задача 1. (7 баллов) Найдите $ f^*(y) $ к следующим функциям}

\subsubsection*{1.1. (0.5 балла)}
$ f : \mathbb{R} \to \mathbb{R} $: 
$$
f(x) = \max \{ |x|, x^2 \}.
$$

\textbf{1.1 решение}

Функция $f(x)$ может быть записана в виде:
$$
f(x) = 
\begin{cases} 
|x|, & \text{если } |x| \leq 1, \\
x^2, & \text{если } |x| \geq 1.
\end{cases}
$$

Рассмотрим функцию  $g(x) = x y - f(x)$ и найдем супремум по $ x \in \mathbb{R} $, разбивая область на две части: $ |x| \leq 1 $ и $|x| \geq 1 $.

\vspace{5 mm}

1. $|x| \leq 1$

\begin{itemize}
\item При  $x \in [0, 1]$:  $g(x) = x y - x = x(y - 1) $.
  \begin{itemize}
  \item Если $ y > 1 $, максимум при $ x = 1 $: $ g(1) = y - 1 $.
  \item Если $ y = 1 $, $ g(x) = 0 $.
  \item Если $ y < 1 $, максимум при $ x = 0 $: $ g(0) = 0 $.
  \end{itemize}
  
\item При $ x \in [-1, 0] $: $ g(x) = x y + x = x(y + 1) $.
  \begin{itemize}
  \item Если $ y > -1 $, максимум при $ x = 0 $: $ g(0) = 0 $.
  \item Если $ y = -1 $, $ g(x) = 0 $.
  \item Если $ y < -1 $, максимум при $ x = -1 $: $ g(-1) = -y - 1 $.
  \end{itemize}
\end{itemize}

Итог для $ |x| \leq 1 $:
$$
\begin{cases}
y - 1, & \text{если } y > 1, \\
0, & \text{если } -1 \leq y \leq 1, \\
-y - 1, & \text{если } y < -1.
\end{cases}
$$

\vspace{15 mm}

2. $ |x| \geq 1 $

\begin{itemize}
\item При $ x \geq 1 $: $ g(x) = x y - x^2 $. 
  \begin{itemize}
  \item $ x = \frac{y}{2} $. Она лежит в $ x \geq 1 $ при $ y \geq 2 $, значение: $ g(\frac{y}{2}) = \frac{y^2}{4} $.
  \item При $ y < 2 $, максимум на $ x \geq 1 $ при $ x = 1 $: $ g(1) = y - 1 $.
  \end{itemize}
  
\item При $ x \leq -1 $: $ g(x) = x y - x^2 $.
  \begin{itemize}
  \item $ x = \frac{y}{2} $. Она лежит в $ x \leq -1 $ при $ y \leq -2 $, значение: $ g(\frac{y}{2}) = \frac{y^2}{4} $.
  \item При $ y > -2 $, максимум на $ x \leq -1 $ при $ x = -1 $: $ g(-1) = -y - 1 $.
  \end{itemize}
\end{itemize}

Итог для $ |x| \geq 1 $:
$$
B(y) = 
\begin{cases}
\frac{y^2}{4}, & \text{если } y \geq 2, \\
y - 1, & \text{если } 0 < y < 2, \\
-y - 1, & \text{если } -2 < y < 0, \\
\frac{y^2}{4}, & \text{если } y \leq -2.
\end{cases}
$$

Расмотрим несколько случаев:

\begin{itemize}
\item При $ y > 2 $: $ A(y) = y - 1 $, $ B(y) = \frac{y^2}{4} $, причем $ \frac{y^2}{4} > y - 1 $, поэтому $ f^*(y) = \frac{y^2}{4} $.
\item При $ 1 < y < 2 $: $ A(y) = y - 1 $, $ B(y) = y - 1 $, поэтому $ f^*(y) = y - 1 $.
\item При $ -1 \leq y \leq 1 $: $ A(y) = 0 $, $ B(y) \leq 0 $, поэтому $ f^*(y) = 0 $.
\item При $ -2 < y < -1 $: $ A(y) = -y - 1 $, $ B(y) = -y - 1 $, поэтому $ f^*(y) = -y - 1 $.
\item При $ y < -2 $: $ A(y) = -y - 1 $, $ B(y) = \frac{y^2}{4} $, причем $ \frac{y^2}{4} > -y - 1 $, поэтому $ f^*(y) = \frac{y^2}{4} $.
\end{itemize}

получается:
$$
f^*(y) = 
\begin{cases} 
0, & \text{если } |y| \leq 1, \\
|y| - 1, & \text{если } 1 < |y| < 2, \\
\frac{y^2}{4}, & \text{если } |y| \geq 2.
\end{cases}
$$

\subsubsection*{1.2. (0.5 балла)}
$ f : \mathbb{R} \to \mathbb{R} $: 
$$
f(x) = (|x| + 1) \log(|x| + 1).
$$
\textit{Указание: можно пользоваться выпуклостью $ f(x) $ без доказательства.}

\textbf{1.2 решение}

заметим что сопряженная функция будет симметричной (так как сама функция симметрична)

Б.О.О. будем находить $f^*(y)$ для $y > 0$

ясно что супремум достигается при $x>=0$ 
$$
\begin{aligned}
    f^*(y)
    &= \sup_{x \in \mathbb{R}} \{ x y - f(x) \} \\
    &= \sup_{x \in \mathbb{R}} \{ x y - (|x| + 1) \log(|x| + 1) \\
    &= \sup_{x \geq 0} \{ x y - (x + 1) \log(x + 1)
\end{aligned}
$$

$$
\begin{aligned}
    g(x) &= x y - (x + 1) \log(x + 1) \\
    g'(x) &= y - \log(x + 1) - 1
\end{aligned}
$$

заметим что $y - \log(x + 1) - 1  =0$ возможно при $y > 1$

так как производная функции f возрастает, то взглянув на производную в 0 можно понять что субградиент $\partial f(x_0) = [-1, 1]$ где $x_0 = 0$ (это значит что функция уже подпирается всеми прямыми $y \in [-1, 1]$ и в этом случае сдвиг не нужен, он равен 0)

при $y < 0$ все аналогично, только используем $|y|$
$$
f^*(y)
\begin{cases}
    e^{y-1} - y, \quad |y| > 1 \\
    0, \quad |y| \leq 1
\end{cases}
$$



\subsubsection*{1.3. (1 балл)}
$ f : \mathbb{R}^d \to \mathbb{R} $: 
$$
f(x) = (\|x\|_2 + 1) \log(\|x\|_2 + 1) - \|x\|_2.
$$
\textit{Указание: можно пользоваться выпуклостью $ f(x) $ без доказательства.}

\textbf{1.2 решение}

$$
\begin{aligned}
    f^*(y)
    &= \sup_{x \in \mathbb{R}} \{ \langle x, y\rangle - f(x) \} \\
    &= \sup_{x \in \mathbb{R}} \{ \langle x, y\rangle - (\|x\|_2 + 1) \log(\|x\|_2 + 1) + \|x\|_2 \\
    &= \sup_{\alpha > 0} \{\sup_{\|x\|_2  =\alpha}\{\langle x, y\rangle - (\alpha + 1) \log(\alpha + 1) + \alpha\}\}  \\
    &= \sup_{\alpha > 0} \{\alpha \|y\|_2- (\alpha + 1) \log(\alpha + 1) + \alpha\} \\
    & \text{супремум находится легко, так что я сразу напишу ответ} \\
    &= e^{\|y\|_2} - \|y\| - 1
\end{aligned}
$$
\subsubsection*{1.4. (1 балл)}
$ f : \mathbb{R}^d \to \mathbb{R} $: 
$$
f(x) = \max_{i=1,d} x_i.
$$

\textbf{1.4 решение}

$$
\begin{aligned}
    f^*(y)
    &= \sup_{x \in \mathbb{R}} \{ \langle x, y\rangle - f(x) \}
\end{aligned}
$$

допустим в y есть отрицательная компонента. Тогда существует x такой что

$x = (0, 0, ..., -t, ..., 0)$  где нулевая соответсвующая компонента где t > 0. То есть $f^*(y) = y_i(-t) - 0 \xrightarrow{t  arrow \infty} \infty$

допусти что $\sum y_i\ne 1$. 

$x = (t, ..., t)$. Тогда $f^*(y) = \sum y_i * t -t = t (\sum y_i - 1)  arrow \infty $
при стремлении t в соответсвующую бесконечность

если $y_i > 0 \quad \text{и} \quad \sum y_i = 1 $

$$
\langle y, x \rangle = \sum_{i=1}^d y_i x_i \leq \sum_{i=1}^d y_i \max_j x_j = \max_j x_j = f(x).
$$
Следовательно:
$$
\langle y, x \rangle - f(x) \leq 0.
$$
Равенство достигается, например, при  $x = 0$ . Таким образом, $ f^*(y) = 0$ .
$$
f^*(y) = 
\begin{cases} 
0, & \text{если } y_i \geq 0 \text{ и } \sum_{i=1}^d y_i = 1, \\
+\infty, & \text{иначе}.
\end{cases}
$$


\subsubsection*{1.5. (1 балл)}
$ f : \mathbb{R} \to \mathbb{R} $: 
$$
f(x) = \sqrt{1 + x^2}.
$$

\textbf{1.5 решение}

$$
\begin{aligned}
    f^*(y)
    &= \sup_{x \in \mathbb{R}} \{ x y - f(x) \} \\
    &= \sup_{x \in \mathbb{R}} \{ x y - \sqrt{1 + x^2} \} \\
    & \text{легко посчитать супремум при |y| < 1} \\
    &= - \sqrt{1 - y^2}
\end{aligned}
$$

если же $y > 0$  то при $ x y - \sqrt{1 + x^2} \xrightarrow{{x  arrow \infty}} \infty $

$$
f^*(y) = 
\begin{cases} 
- \sqrt{1 - y^2}, & \text{если } |y| < 1 \\
+\infty, & \text{иначе}.
\end{cases}
$$


\subsubsection*{1.6. (1 балл)}
$ f : \mathbb{R} \to \mathbb{R} $: 
$$
f(x) = ch x - 1.
$$

\textbf{1.6 решение}

% -------------------

$$
\begin{aligned}
    f^*(y)
    &= \sup_{x \in \mathbb{R}} \{ x y - f(x) \} \\
    &= \sup_{x \in \mathbb{R}} \{x y - ( ch x - 1) \}
\end{aligned}
$$

$$
g(x) = x y - (ch x - 1) = x y - ch x + 1.
$$

$$
g'(x) = y - sh x.
$$

$$
x = arsh y
$$

$$
ch x = \sqrt{sh^2 x + 1} = \sqrt{y^2 + 1}.
$$

Тогда:
$$
\begin{aligned}
g(x) &= x y - ch x + 1 \\
&= y \cdot arsh y - \sqrt{y^2 + 1} + 1.
\end{aligned}
$$

Получается
$$
f^*(y) = y \cdot arsh y - \sqrt{y^2 + 1} + 1.
$$


\subsubsection*{1.7. (2 балла)}
$ f : \mathbb{S}_{++}^d \to \mathbb{R} $: 
$$
f(X) = Tr(X^{-1}).
$$
\textit{Указание: возможно пригодится, что $ (I_d + tuu^\top)^{-1} = I_d - \frac{tuu^\top}{1+t} $, где $ \|u\|_2 = 1 $.}

\textbf{1.7 решение}

Рассмотрим функцию $g(X) = Tr(Y X) - Tr(X^{-1})$.
$$
d g = Tr(Y^TdX) + Tr(X^{-1}dX X^{-1}) = Tr( (Y + X^{-2}) \, dX ).
$$
Для того чтобы это было равно нулю для всех $dX$, необходимо:
$$
Y + X^{-2} = 0 \quad  arrow \quad X^{-2} = -Y \quad  arrow \quad X = (-Y)^{-1/2}.
$$
Это возможно только если Y строго отрицательно определена. И тогда
$$
g(X) = -Tr(X^{-1}) - Tr(X^{-1}) = -2 Tr(X^{-1}) = -2 Tr((-Y)^{1/2}).
$$

Если $Y$ не отрицательно полуопределена, то существует направление, в котором $Y$ положительно определена, и тогда можно выбрать $X$ таким, что $Tr(Y X)$ неограниченно возрастает, а $Tr(X^{-1})$ ограничено, поэтому супремум равен $+\infty$.
Получаем 

$$
f^*(Y) = 
\begin{cases} 
-2 Tr ( (-Y)^{1/2}  ), & \text{если } Y < 0, \\
+\infty, & \text{иначе}.
\end{cases}
$$

\subsection*{Задача 2. (1 балл)}
Выразите сопряжённую функцию $ f^*(y) $ через $ g^*(y) $ и $ h^*(y) $ для $ f : \mathbb{R}^d \to \mathbb{R} $: 
$$
f(x) = \inf_{u+v=x} (g(u) + h(v)),
$$
где $ g : \mathbb{R}^d \to \mathbb{R} $ и $ h : \mathbb{R}^d \to \mathbb{R} $.

\textbf{2 решение}

% ----------------------------
$$
\begin{aligned}
    f^*(y)
    &= \sup_{x \in \mathbb{R}^d} \{ \langle y, x \rangle - f(x) \} \\
    &= \sup_{x \in \mathbb{R}^d} \{ \langle y, x \rangle - \inf_{u+v=x} (g(u) + h(v)) \} \\
    &= \sup_{x \in \mathbb{R}^d} \{  \sup_{u+v=x} (\langle y, x \rangle - g(u) - h(v)) \} \\
    &= \sup_{x \in \mathbb{R}^d} \{  \sup_{u+v=x} (\langle y, u \rangle - g(u) + \langle y, v \rangle- h(v)) \} \\
    &= g^*(y) + h^*(y)
\end{aligned}
$$

\subsection*{Задача 3. (2 балла)}
Найдите $ f^{**}(x) $ для $ f : \mathbb{R}^d \to \mathbb{R} $:
$$
f(x) = \mathbb{I}\{x \in D\} = 
\begin{cases} 
0 & x \in D, \\ 
+\infty & x \notin D,
\end{cases}
$$
где $ D \subset \mathbb{R}^d $ компактно.

\textbf{3 решение}


$$
f^*(y) = \sup_{x \in \mathbb{R}^d}  \{ \langle y, x \rangle - f(x)  \} = \sup_{x \in D} \langle y, x \rangle,
$$

Тогда:
$$
f^{**}(x) = \sup_{y \in \mathbb{R}^d}  \{ \langle x, y \rangle - f^*(y)  \} = \sup_{y \in \mathbb{R}^d}  \{ \langle x, y \rangle - \sup_{z \in D} \langle y, z \rangle  \} = \sup_{y \in \mathbb{R}^d} \inf_{z \in D} \langle y, x - z \rangle.
$$

Рассмотрим два случая:


Если $ x \in conv(D) $, то для любого $ y \in \mathbb{R}^d $ выполняется:
$$
\langle x, y \rangle \leq \sup_{z \in D} \langle z, y \rangle = f^*(y),
$$
следовательно:
$$
\langle x, y \rangle - f^*(y) \leq 0.
$$
При $ y = 0 $ достигается равенство $ 0 $, поэтому $ f^{**}(x) = 0 $.

Если $ x \notin conv(D) $, то по теореме отделимости существует вектор $ y \in \mathbb{R}^d $ и число $ \alpha \in \mathbb{R} $ такие, что:
$$
\langle x, y \rangle > \alpha \quad \text{и} \quad \langle z, y \rangle < \alpha \quad \text{для всех } z \in D.
$$
Тогда:
$$
\langle x, y \rangle - f^*(y) = \langle x, y \rangle - \sup_{z \in D} \langle z, y \rangle \geq \langle x, y \rangle - \alpha > 0.
$$
Рассматривая $ t y $ при $ t \to +\infty $, получаем:
$$
\langle x, t y \rangle - f^*(t y) = t  ( \langle x, y \rangle - f^*(y)  ) \to +\infty,
$$
поэтому $ f^{**}(x) = +\infty $.


Таким образом, второе сопряжение $ f^{**}(x) $ является индикаторной функцией выпуклой оболочки множества $ D $:
$$
f^{**}(x) = \mathbb{I}\{x \in conv(D)\}.
$$


$$
f(x) = 
\begin{cases} 
0, & \text{если } x \in conv(D), \\
+\infty, & \text{если } x \notin conv(D).
\end{cases}
$$

\end{document}
